% Auto-generated by pipeline/scripts/amplicon_exon_overlap_table.py
% via rule amplicon_exon_overlap_table
\begin{table}[H]
\centering
\small
\begin{tabular}{clrrrrrr}
\toprule
\textbf{Exon} & \textbf{Amplicon} & \textbf{Amp.~start} & \textbf{Amp.~end} & \textbf{Amp.~len} & \textbf{Exon start} & \textbf{Exon end} & \textbf{Exon len} \\
 & & & & \textbf{(bp)} & & & \textbf{(bp)} \\
\midrule
18 & Amplicon~2 & 55{,}241{,}583 & 55{,}241{,}796 & 213 & 55{,}241{,}613 & 55{,}241{,}736 & 123 \\
19 & Amplicon~1 & 55{,}242{,}378 & 55{,}242{,}574 & 196 & 55{,}242{,}414 & 55{,}242{,}513 & 99 \\
20 & Amplicon~4 & 55{,}248{,}956 & 55{,}249{,}194 & 238 & 55{,}248{,}985 & 55{,}249{,}171 & 186 \\
21 & Amplicon~3 & 55{,}259{,}374 & 55{,}259{,}589 & 215 & 55{,}259{,}411 & 55{,}259{,}567 & 156 \\
\bottomrule
\end{tabular}
\caption{Amplicon spans relative to their target EGFR exon (NM\_005228, hg19, plus strand, 0-based half-open coordinates). Each amplicon fully contains its target exon; the primers are placed in the flanking intronic sequence on either side. Left intronic flank = exon~start $-$ amplicon~start; right intronic flank = amplicon~end $-$ exon~end.}
\label{tab:amplicon-exon-overlap}
\end{table}
